\section{Introduction}

A developmental phenotype is a trajectory through interacting constraints. Genes, tissue mechanics, endocrine signaling, immune activity, experience, and available environments can affect one another over time. A diagnosis may summarize recurring outcomes without identifying the particular sequence that produced them. The practical challenge is to describe this coupling without converting every plausible association into a causal explanation.

The present synthesis makes one organizing proposal: constitutional variation and early regulatory conditions may converge partly through synaptic developmental machinery, altering the conditions under which sensory processing, autonomic regulation, attention, affiliation, and action become coordinated. The central sequence is Constitution \ensuremath{\rightarrow} Perinatal Initialization \ensuremath{\rightarrow} \ensuremath{\alpha_2\delta}–CaV2 Hinge \ensuremath{\rightarrow} Network Calibration \ensuremath{\rightarrow} Regulatory Phenotype \ensuremath{\rightarrow} Developmental Trajectory. Its arrows have different evidentiary status. The sequence is a research scaffold, not an established biological pathway.

The broader environmental theory asks a separate question: when do institutions, tools, and human–AI coupling change the functional payoff of a regulatory configuration? Its validity does not depend on RCCX or perinatal molecular mediation. This separation makes the architecture modular: one bridge can fail without immunizing, or automatically invalidating, the remaining claims.


\subsection{Corpus, synthesis method, and evidence grammar}

The canonical project corpus comprises the Cascade v3.2, The Polymathic Singularity v0.1, and the Perinatal Regulatory Initialization / \ensuremath{\alpha_2\delta}–CaV2 extension \cite{projP1,projP2,projP3}. This article resolves overlaps by using Regulatory Coupling Orientation for the qualitative construct and Regulatory Closure Bias (RCB) for its proposed quantitative index. Perinatal Regulatory Initialization (PRI) names an exposure-and-response interval, not a new organ or a single biomarker. “Hinge” denotes a candidate convergence mechanism, not an obligatory bottleneck.

Targeted literature searches on 25 September 2026 examined the supplied anchors, primary molecular studies, human cohort and intervention findings, and counterevidence. Searches included birth signaling with oxytocin-null results; \emph{L. reuteri} with restricted or null outcomes; birth mode with adjusted null associations; channel-subtype switching with synapse-specific constraints; and genotype–phenotype evidence at the RCCX interface. This is a critical narrative synthesis. It is neither a systematic review nor a quantitative meta-analysis, and no new participants or experimental observations are reported.

Four labels govern claims. Established mechanism means experimental support for the specified mechanism in the specified preparation. Supported synthesis means convergent or directly relevant evidence with remaining inferential distance. Proposed mechanism means a testable connection not directly demonstrated. Horizon hypothesis means a broader explanatory interpretation whose constructs or causal links require development. An established mouse mechanism is not thereby an established human mechanism.

Design codes add resolution: E0 conceptual; E1 observational; E2 replicated observational; E3 experimental animal or cell work; E4 human intervention or quasi-experiment; E5 replicated human causal evidence. These codes classify evidence types, not probabilities of truth or a universal quality ranking. Randomization can identify an intervention effect while leaving its proposed mediator unmeasured. No claim in this manuscript receives E5 for the complete cascade.


\subsection{Prior frameworks and division of explanatory labor}

The Cascade joins constitutional sensitivity, state-dependent regulation, developmental reorganization, and conversion of capacity into output \cite{projP1}. Its useful contribution is recursive organization. Attention selects environments; environments alter load; physiological state affects attention; social response changes subsequent opportunities. These relations must be measured separately rather than compressed into a diagnostic identity.

The PRI extension supplies an upstream developmental interval and a candidate molecular interface \cite{projP3}. The Polymathic Singularity concerns the changing ecological value of already developing or developed regulatory configurations \cite{projP2}. It proposes that cognitive externalization and AI-assisted execution can reduce translation costs between ideas and artifacts. It also predicts branch proliferation, unfinished work, and destabilization when amplification exceeds closure. These are environmental interaction hypotheses, not evidence for a genetic syndrome.

The three-toggle account is retained as a conceptual decomposition: available integrative capacity, direction of reorganization, and conversion/dissemination. The product metaphor is not a measured equation of historical achievement. “Coherence” must be operationalized through functioning, goal consistency, recovery, and completed work; it cannot mean agreement with the theory or moral superiority.
